% ============================================================================
% APPENDIX FI: LIT-FISCHBACHER - Conditional Cooperation & Experimental Methods
% Evidence-Based Framework for Economic and Social Behavior
% ============================================================================
% Category: LIT- (Literature Integration)
% Language: English
% Code: FI
% Version: 1.0
% Last Updated: 2026-01-17
% Primary Dependencies: CORE-HOW, METHOD-E, LIT-FEHR, LIT-GAECHTER
% EBF Integration: Conditional cooperation types and experimental methodology
% ============================================================================

\chapter*{Appendix UNMAPPED_FI: Urs Fischbacher Research Integration}
\addcontentsline{toc}{chapter}{Appendix UNMAPPED_FI: LIT-FISCHBACHER - Conditional Cooperation \& Experimental Methods}
\label{app:fischbacher}

% ============================================================================
% HEADER BLOCK
% ============================================================================
\begin{tcolorbox}[colback=white,colframe=black!75!white,title=Appendix Metadata]
\textbf{Appendix:} FI -- LIT-FISCHBACHER \\
\textbf{Category:} Literature Integration (LIT-) \\
\textbf{Version:} 1.0 (January 2026) \\
\textbf{Dependencies:} CORE-HOW, METHOD-E, LIT-FEHR, LIT-GAECHTER \\
\textbf{Status:} Complete
\end{tcolorbox}

% ============================================================================
% CROSS-REFERENCE MAP
% ============================================================================
\begin{tcolorbox}[colback=blue!5!white,colframe=blue!75!black,title=Cross-Reference Map]
\textbf{Builds On:}
\begin{itemize}
    \item Appendix K (LIT-FEHR): Social preferences foundations
    \item Appendix AM (LIT-GAECHTER): Cooperation and punishment
    \item Appendix UNMAPPED_B (CORE-HOW): Complementarity in social interaction
\end{itemize}

\textbf{Supports:}
\begin{itemize}
    \item Appendix E (METHOD-OPS): Experimental operationalization
    \item Appendix AJ (DOMAIN-SOCIAL): Social preferences applications
    \item Appendix UNMAPPED_AU (CORE-AWARE): Awareness in cooperation decisions
\end{itemize}

\textbf{Related Literature:}
\begin{itemize}
    \item Appendix UNMAPPED_CA (LIT-CAMERER): Strategic behavior in games
    \item Appendix W (LIT-ARIELY): Dishonesty research
    \item Appendix UNMAPPED_SU (LIT-SUTTER): Team decision making
\end{itemize}
\end{tcolorbox}

% ============================================================================
% CHAPTER LINKAGE
% ============================================================================
\begin{tcolorbox}[colback=green!5!white,colframe=green!75!black,title=Chapter Linkage]
\textbf{Primary:} Chapter 5 (Complementarity) -- Social preference types and interaction effects

\textbf{Secondary:}
\begin{itemize}
    \item Chapter 9 (Context): Beliefs about others as context
    \item Chapter 10.1 (Social Dimension): FEPSDE social preferences
    \item Chapter 14 (Applications): Public goods and cooperation
\end{itemize}
\end{tcolorbox}

% ============================================================================
% ABSTRACT
% ============================================================================
\begin{abstract}
This appendix integrates Urs Fischbacher's research on conditional cooperation, experimental methodology, and social norms into the Evidence-Based Framework. Fischbacher's work establishes that cooperation types are heterogeneous and measurable: approximately 50\% are conditional cooperators, 30\% free-riders, and 14\% unconditional cooperators. His creation of z-Tree revolutionized experimental economics methodology. We formalize six core axioms (UNMAPPED_FI-1 to UNMAPPED_FI-6) connecting cooperation dynamics to EBF. Key parameters include conditional cooperation slope ($\beta_{CC} \approx 0.5$), cooperation decay rate, and partial lying prevalence. Co-author with Ernst Fehr on foundational social preferences papers.
\end{abstract}

% ============================================================================
% QUICK REFERENCE
% ============================================================================
\begin{tcolorbox}[colback=gray!5!white,colframe=gray!75!black,title=Quick Reference]
\textbf{Researcher:} Urs Fischbacher (University of Konstanz)

\textbf{Core Contributions:}
\begin{itemize}
    \item Conditional cooperation typology and measurement
    \item z-Tree experimental software (8,000+ citations)
    \item Cooperation decay mechanisms
    \item Partial dishonesty and self-image
\end{itemize}

\textbf{Key Parameters:}
\begin{itemize}
    \item Conditional cooperators: $\approx 50\%$ of population
    \item Free-riders: $\approx 30\%$ of population
    \item Unconditional cooperators: $\approx 14\%$ of population
    \item Conditional cooperation slope: $\beta_{CC} \approx 0.5$
    \item Partial lying: Most cheaters don't maximize
\end{itemize}

\textbf{EBF Parameters Updated:}
\begin{itemize}
    \item Cooperation type distribution $\pi = (\pi_{CC}, \pi_{FR}, \pi_{UC})$
    \item Conditional slope $\beta_{CC}$ in CORE-HOW
    \item Belief-driven dynamics in cooperation
\end{itemize}
\end{tcolorbox}

% ============================================================================
% FUNDAMENTAL QUESTION
% ============================================================================
% -----------------------------------------------------------------------------
% APPENDIX SCOPE BOX
% -----------------------------------------------------------------------------

\begin{tcolorbox}[
    colback=orange!5!white,
    colframe=orange!75!black,
    title={\textbf{Appendix Scope: Ziel / In-Scope / Out-of-Scope / Constraints}},
    fonttitle=\bfseries
]

\textbf{Ziel (Objective):}
\begin{itemize}[nosep]
    \item Why do some people cooperate while others free-ride, and what mechanisms explain the dynamics of cooperation in groups?
\end{itemize}

\textbf{In-Scope:}
\begin{itemize}[nosep]
    \item Fundamental Question
    \item Core Axioms: Fischbacher's Contributions
    \item Key Empirical Results
    \item Integration with Evidence-Based Framework
\end{itemize}

\textbf{Out-of-Scope (delegiert an andere Appendices):}
\begin{itemize}[nosep]
    \item LIT-FEHR $\rightarrow$ Appendix K
    \item LIT-GAECHTER $\rightarrow$ Appendix AM
    \item CORE-HOW $\rightarrow$ Appendix UNMAPPED_B
    \item METHOD-OPS $\rightarrow$ Appendix E
    \item DOMAIN-SOCIAL $\rightarrow$ Appendix AJ
\end{itemize}

\textbf{Constraints (Anwendungsgrenzen):}
\begin{itemize}[nosep]
    \item [TODO: Define when this appendix does NOT apply]
    \item [TODO: Define required prerequisites]
    \item [TODO: Define known limitations]
\end{itemize}

\textbf{Lieferobjekte (Deliverables):}
\begin{itemize}[nosep]
    \item 2 Table(s)
    \item 9 Equation(s)
    \item 12 Axiom(s)
    \item Worked Example(s)
\end{itemize}

\end{tcolorbox}


\section{Fundamental Question}

\begin{tcolorbox}[colback=red!5!white,colframe=red!75!black,title=Central Question]
\textbf{Why do some people cooperate while others free-ride, and what mechanisms explain the dynamics of cooperation in groups?}
\end{tcolorbox}

Fischbacher's research addresses cooperation puzzles:

\begin{enumerate}
    \item \textbf{Heterogeneity}: Not everyone has the same social preferences
    \item \textbf{Conditionality}: Most cooperators condition on others' behavior
    \item \textbf{Dynamics}: Cooperation often decays over time---why?
    \item \textbf{Measurement}: How to reliably identify cooperation types?
\end{enumerate}

For EBF, Fischbacher provides the empirical foundation for understanding how complementarities in social behavior ($\gamma > 0$) operate through heterogeneous types.

% ============================================================================
% THEORETICAL FOUNDATIONS
% ============================================================================
\section{Theoretical Foundations}

\subsection{Cooperation Type Taxonomy}

The strategy method reveals distinct behavioral types:

\begin{equation}
g_i(g_{-i}) = \begin{cases}
\beta_{CC} \cdot \bar{g}_{-i} & \text{Conditional Cooperator (50\%)} \\
0 & \text{Free-Rider (30\%)} \\
g_{max} & \text{Unconditional Cooperator (14\%)} \\
\text{other} & \text{Triangle/Hump-shaped (6\%)}
\end{cases}
\end{equation}

where $g_i$ = individual contribution, $\bar{g}_{-i}$ = average others' contribution.

\subsection{Cooperation Decay Mechanism}

Why does cooperation decay even with conditional cooperators?

\begin{equation}
g_t = \beta_{CC} \cdot \mathbb{E}_t[\bar{g}_{-i}] + \epsilon_t \quad \text{with} \quad \mathbb{E}_t[\bar{g}_{-i}] < g_{t-1}
\end{equation}

\textbf{Key insight}: Imperfect conditional cooperation plus pessimistic belief updating drives decay. Conditional cooperators contribute \textit{slightly less} than others' average, creating downward spiral.

\subsection{The z-Tree Revolution}

Fischbacher's z-Tree software (2007) standardized experimental economics:

\begin{itemize}
    \item Enabled complex interactive experiments
    \item Improved replicability across labs
    \item 8,000+ citations, used in thousands of studies
    \item Foundation of modern experimental economics infrastructure
\end{itemize}

% ============================================================================
% CORE AXIOMS
% ============================================================================
\section{Core Axioms: Fischbacher's Contributions}

\subsection{Axiom UNMAPPED_FI-1: Cooperation Type Heterogeneity}

\begin{tcolorbox}[colback=blue!5!white,colframe=blue!75!black,title=Axiom UNMAPPED_FI-1: Cooperation Type Heterogeneity]
\textbf{Statement}: The population consists of distinct cooperation types with stable, measurable preferences. Approximately 50\% are conditional cooperators, 30\% free-riders, 14\% unconditional cooperators.

\textbf{Formal Definition}:
\begin{equation}
\pi = (\pi_{CC}, \pi_{FR}, \pi_{UC}, \pi_{other}) \approx (0.50, 0.30, 0.14, 0.06)
\end{equation}

\textbf{Empirical Support}:
\begin{itemize}
    \item Fischbacher, Gächter, Fehr (2001): Original classification
    \item Replicated across 15+ countries
    \item Stable within individuals over time
\end{itemize}

\textbf{Evidence Tier}: 1 (Replicated experimental finding)
\end{tcolorbox}

\subsection{Axiom UNMAPPED_FI-2: Conditional Cooperation Function}

\begin{tcolorbox}[colback=blue!5!white,colframe=blue!75!black,title=Axiom UNMAPPED_FI-2: Conditional Cooperation Function]
\textbf{Statement}: Conditional cooperators' contributions are a positive function of others' contributions, with slope $\beta_{CC} \approx 0.5$.

\textbf{Formal Definition}:
\begin{equation}
g_i = \alpha + \beta_{CC} \cdot \bar{g}_{-i} + \epsilon, \quad \beta_{CC} \in [0.4, 0.6]
\end{equation}

\textbf{Key Feature}: $\beta_{CC} < 1$ means imperfect matching---conditional cooperators contribute slightly less than the average, creating decay.

\textbf{Evidence Tier}: 1 (Robust experimental finding)
\end{tcolorbox}

\subsection{Axiom UNMAPPED_FI-3: Belief-Driven Dynamics}

\begin{tcolorbox}[colback=blue!5!white,colframe=blue!75!black,title=Axiom UNMAPPED_FI-3: Belief-Driven Cooperation Decay]
\textbf{Statement}: Cooperation decay is primarily driven by belief updating, not preference change. Conditional cooperators update beliefs pessimistically.

\textbf{Formal Definition}:
\begin{equation}
\mathbb{E}_t[g_{-i}] = (1-\lambda) \cdot \mathbb{E}_{t-1}[g_{-i}] + \lambda \cdot g_{-i,t-1}
\end{equation}

with observed $g_{-i,t-1} < \mathbb{E}_{t-1}[g_{-i}]$ triggering downward revision.

\textbf{Empirical Support}:
\begin{itemize}
    \item Fischbacher \& Gächter (2010): Beliefs explain decay better than preference change
    \item Direct belief elicitation confirms pessimistic updating
\end{itemize}

\textbf{Evidence Tier}: 1 (Causal identification through belief elicitation)
\end{tcolorbox}

\subsection{Axiom UNMAPPED_FI-4: Third-Party Punishment}

\begin{tcolorbox}[colback=blue!5!white,colframe=blue!75!black,title=Axiom UNMAPPED_FI-4: Third-Party Punishment]
\textbf{Statement}: Uninvolved third parties punish norm violators at personal cost, extending norm enforcement beyond direct interactions.

\textbf{Formal Definition}:
\begin{equation}
P(\text{punish}|\text{norm violation observed}) \approx 0.6
\end{equation}

even when punisher is not directly affected.

\textbf{Implications}:
\begin{itemize}
    \item Enables large-scale cooperation
    \item Creates decentralized norm enforcement
    \item $\gamma_{punishment, cooperation} > 0$
\end{itemize}

\textbf{Evidence Tier}: 1 (Fehr \& Fischbacher, 2004)
\end{tcolorbox}

\subsection{Axiom UNMAPPED_FI-5: Partial Dishonesty}

\begin{tcolorbox}[colback=red!5!white,colframe=red!75!black,title=Axiom UNMAPPED_FI-5: Partial Dishonesty]
\textbf{Statement}: When given opportunity to cheat undetected, most people cheat partially rather than maximally. Self-image concerns limit dishonesty.

\textbf{Formal Definition}:
\begin{equation}
\text{Lie}_i = \arg\max \{ \pi(lie) - c_{self-image}(lie) \}
\end{equation}

where $c_{self-image}$ is convex, leading to interior solutions.

\textbf{Empirical Support}:
\begin{itemize}
    \item Die-rolling paradigm: Distribution shifted but not to maximum
    \item Consistent with Ariely's ``fudge factor'' concept
\end{itemize}

\textbf{Evidence Tier}: 1 (Fischbacher \& Föllmi-Heusi, 2013)
\end{tcolorbox}

\subsection{Axiom UNMAPPED_FI-6: Strategy Method Validity}

\begin{tcolorbox}[colback=green!5!white,colframe=green!75!black,title=Axiom UNMAPPED_FI-6: Strategy Method Validity]
\textbf{Statement}: The strategy method (eliciting complete contingent strategies) produces behaviorally equivalent results to direct-response methods.

\textbf{Formal Definition}:
\begin{equation}
g^{strategy}(\bar{g}_{-i}) \approx g^{direct}|\bar{g}_{-i}
\end{equation}

\textbf{Implications}:
\begin{itemize}
    \item Validates cooperation type measurement
    \item Enables richer experimental designs
    \item Foundation for conditional cooperation research
\end{itemize}

\textbf{Evidence Tier}: 1 (Methodological validation)
\end{tcolorbox}

% ============================================================================
% KEY EMPIRICAL RESULTS
% ============================================================================
\section{Key Empirical Results}

\subsection{Parameter Estimates}

\begin{table}[h]
\centering
\caption{Key Parameters from Fischbacher's Research}
\label{tab:fischbacher-parameters}
\begin{tabular}{llcc}
\toprule
\textbf{Parameter} & \textbf{Context} & \textbf{Estimate} & \textbf{Source} \\
\midrule
$\pi_{CC}$ & Conditional cooperators & 50\% & FGF 2001 \\
$\pi_{FR}$ & Free-riders & 30\% & FGF 2001 \\
$\pi_{UC}$ & Unconditional cooperators & 14\% & FGF 2001 \\
$\beta_{CC}$ & Cooperation slope & 0.5 & FGF 2001 \\
Third-party punishment & Punishment rate & 60\% & FF 2004 \\
Partial lying & Cheaters who maximize & $< 20\%$ & FFH 2013 \\
\bottomrule
\end{tabular}
\end{table}

\subsection{Worked Example: Public Goods Decay}

\begin{tcolorbox}[colback=green!5!white,colframe=green!75!black,title=Worked Example: Why Cooperation Decays]
\textbf{Setup}: 10-round public goods game, 4 players, MPCR = 0.4

\textbf{Round 1}: All contribute 50\% of endowment

\textbf{Mechanism}:
\begin{enumerate}
    \item Conditional cooperators observe average = 50\%
    \item With $\beta_{CC} = 0.5$, they contribute $0.5 \times 50\% = 25\%$ plus noise
    \item Some contribute 30\%, some 45\%---average drops to 40\%
    \item Next round: CC's target $0.5 \times 40\% = 20\%$
    \item Spiral continues downward
\end{enumerate}

\textbf{Result}: Cooperation decays from 50\% to $\sim$10\% by round 10

\textbf{EBF Interpretation}: Belief updating ($\Psi_{beliefs}$) combined with imperfect conditional cooperation ($\beta_{CC} < 1$) creates endogenous decay.
\end{tcolorbox}

% ============================================================================
% EBF INTEGRATION
% ============================================================================
\section{Integration with Evidence-Based Framework}

\subsection{Connection to 10C Framework}

\begin{table}[h]
\centering
\caption{Fischbacher's Contributions to 10C Framework}
\label{tab:fischbacher-9c}
\begin{tabular}{lll}
\toprule
\textbf{CORE} & \textbf{Fischbacher Contribution} & \textbf{Mechanism} \\
\midrule
CORE-WHO (AAA) & Type distribution & $\pi = (\pi_{CC}, \pi_{FR}, \pi_{UC})$ \\
CORE-HOW (B) & Conditional complementarity & $\gamma_{ij} > 0$ for CC types \\
CORE-WHEN (V) & Beliefs as context & $\Psi_{beliefs}$ drives dynamics \\
CORE-AWARE (AU) & Self-image in honesty & $c_{self-image}$ limits cheating \\
\bottomrule
\end{tabular}
\end{table}

\subsection{Cooperation Type Distribution as EBF Parameter}

The cooperation type distribution becomes a fundamental EBF parameter:

\begin{equation}
\bar{g}(\Psi) = \pi_{CC} \cdot g_{CC}(\Psi) + \pi_{FR} \cdot 0 + \pi_{UC} \cdot g_{max}
\end{equation}

This allows prediction of aggregate cooperation given context $\Psi$.

% ============================================================================
% CRITICAL FOUNDATIONS
% ============================================================================
\section{Critical Foundations}

\begin{tcolorbox}[colback=orange!5!white,colframe=orange!75!black,title=Critical Foundations]
\begin{enumerate}
    \item \textbf{Fehr \& Schmidt (1999)}: Inequity aversion model
    \item \textbf{Bolton \& Ockenfels (2000)}: ERC model
    \item \textbf{Selten (1967)}: Strategy method origin
    \item \textbf{Isaac, Walker, Thomas (1984)}: Public goods experiments
    \item \textbf{Marwell \& Ames (1981)}: Free-rider problem experiments
\end{enumerate}
\end{tcolorbox}

% ============================================================================
% OPEN ISSUES
% ============================================================================
\section{Open Issues and Research Frontier}

\begin{enumerate}
    \item \textbf{Type Stability}: How stable are cooperation types across domains?
    \item \textbf{Type Origins}: What determines whether someone is CC vs. FR?
    \item \textbf{Intervention Design}: How to shift type distribution?
    \item \textbf{Large Groups}: Does conditional cooperation scale?
    \item \textbf{Online Behavior}: Cooperation types in digital environments?
\end{enumerate}

% ============================================================================
% SUMMARY
% ============================================================================

% === AUTO-GENERATED ENTRIES (2026-02-22) ===

%%%%%%%%%%%%%%%%%%%%%%%%%%%%%%%%%%%%%%%%%%%%%%%%%%%%%%%%%%%%%%%%%%%%%%%%%%%%%%%
\subsubsection{2006: Heterogeneous Social Preferences and the Dynamics of Free Ri...}
%%%%%%%%%%%%%%%%%%%%%%%%%%%%%%%%%%%%%%%%%%%%%%%%%%%%%%%%%%%%%%%%%%%%%%%%%%%%%%%

\paragraph{Citation.}
Fischbacher, Urs and Gächter, Simon (2006). ``Heterogeneous Social Preferences and the Dynamics of Free Riding in Public Goods.'' \textit{CESifo Working Paper}.

\paragraph{10C Integration.}
\begin{itemize}[nosep]
  \item \textbf{Domain}: [To be specified]
  \item \textbf{use\_for}: LIT-FISCHBACHER, LIT-O
\end{itemize}


%%%%%%%%%%%%%%%%%%%%%%%%%%%%%%%%%%%%%%%%%%%%%%%%%%%%%%%%%%%%%%%%%%%%%%%%%%%%%%%
\subsubsection{2007: z-Tree: Zurich Toolbox for Ready-made Economic Experiments}
%%%%%%%%%%%%%%%%%%%%%%%%%%%%%%%%%%%%%%%%%%%%%%%%%%%%%%%%%%%%%%%%%%%%%%%%%%%%%%%

\paragraph{Citation.}
Fischbacher, Urs (2007). ``z-Tree: Zurich Toolbox for Ready-made Economic Experiments.'' \textit{Experimental Economics}, 10.

\paragraph{10C Integration.}
\begin{itemize}[nosep]
  \item \textbf{Domain}: [To be specified]
  \item \textbf{use\_for}: LIT-FISCHBACHER, LIT-O
\end{itemize}


%%%%%%%%%%%%%%%%%%%%%%%%%%%%%%%%%%%%%%%%%%%%%%%%%%%%%%%%%%%%%%%%%%%%%%%%%%%%%%%
\subsubsection{2013: Lies in Disguise: An Experimental Study on Cheating}
%%%%%%%%%%%%%%%%%%%%%%%%%%%%%%%%%%%%%%%%%%%%%%%%%%%%%%%%%%%%%%%%%%%%%%%%%%%%%%%

\paragraph{Citation.}
Fischbacher, Urs and Föllmi-Heusi, Franziska (2013). ``Lies in Disguise: An Experimental Study on Cheating.'' \textit{Journal of the European Economic Association}, 11.

\paragraph{10C Integration.}
\begin{itemize}[nosep]
  \item \textbf{Domain}: [To be specified]
  \item \textbf{use\_for}: LIT-FISCHBACHER, LIT-O
\end{itemize}


%%%%%%%%%%%%%%%%%%%%%%%%%%%%%%%%%%%%%%%%%%%%%%%%%%%%%%%%%%%%%%%%%%%%%%%%%%%%%%%
\subsubsection{2012: The Behavioral Validity of the Strategy Method in Public Goo...}
%%%%%%%%%%%%%%%%%%%%%%%%%%%%%%%%%%%%%%%%%%%%%%%%%%%%%%%%%%%%%%%%%%%%%%%%%%%%%%%

\paragraph{Citation.}
Fischbacher, Urs and Gächter, Simon and Quercia, Simone (2012). ``The Behavioral Validity of the Strategy Method in Public Good Experiments.'' \textit{Journal of Economic Psychology}, 33.

\paragraph{10C Integration.}
\begin{itemize}[nosep]
  \item \textbf{Domain}: [To be specified]
  \item \textbf{use\_for}: LIT-FISCHBACHER, LIT-O
\end{itemize}


%%%%%%%%%%%%%%%%%%%%%%%%%%%%%%%%%%%%%%%%%%%%%%%%%%%%%%%%%%%%%%%%%%%%%%%%%%%%%%%
\subsubsection{2014: Heterogeneous Reactions to Heterogeneity in Returns from Pub...}
%%%%%%%%%%%%%%%%%%%%%%%%%%%%%%%%%%%%%%%%%%%%%%%%%%%%%%%%%%%%%%%%%%%%%%%%%%%%%%%

\paragraph{Citation.}
Fischbacher, Urs and Schudy, Simeon and Teyssier, Sabrina (2014). ``Heterogeneous Reactions to Heterogeneity in Returns from Public Goods.'' \textit{Social Choice and Welfare}, 43.

\paragraph{10C Integration.}
\begin{itemize}[nosep]
  \item \textbf{Domain}: [To be specified]
  \item \textbf{use\_for}: LIT-FISCHBACHER, LIT-O
\end{itemize}

% === END AUTO-GENERATED ENTRIES ===
\section{Summary}

\begin{tcolorbox}[colback=gray!5!white,colframe=gray!75!black,title=Chapter Summary]
Urs Fischbacher's research provides EBF with:

\textbf{Theoretical Contributions}:
\begin{itemize}
    \item Cooperation type taxonomy (CC, FR, UC)
    \item Belief-driven cooperation dynamics
    \item Third-party punishment mechanism
    \item Partial dishonesty model
\end{itemize}

\textbf{Methodological Contributions}:
\begin{itemize}
    \item z-Tree experimental software
    \item Strategy method validation
    \item Cooperation type measurement
\end{itemize}

\textbf{Key Parameters}:
\begin{itemize}
    \item Type distribution: $\pi = (0.50, 0.30, 0.14, 0.06)$
    \item Conditional slope: $\beta_{CC} \approx 0.5$
    \item Third-party punishment: 60\%
\end{itemize}

\textbf{EBF Integration}:
\begin{itemize}
    \item Population heterogeneity in CORE-WHO
    \item Conditional complementarity in CORE-HOW
    \item Belief dynamics in CORE-WHEN
\end{itemize}

\textbf{Key Insight}: Cooperation dynamics emerge from the interaction of heterogeneous types, with beliefs mediating behavior. Design must account for type distribution.
\end{tcolorbox}

% ============================================================================
% GLOSSARY
% ============================================================================
\section{Glossary}

Key terms defined in this appendix (see also Appendix G for complete Glossary):

\begin{description}
    \item[Conditional Cooperator] Person who contributes proportionally to others' contributions
    \item[Free-Rider] Person who contributes nothing regardless of others
    \item[Strategy Method] Eliciting complete contingent strategies before outcomes realized
    \item[Third-Party Punishment] Punishment by uninvolved observers
    \item[z-Tree] Zurich Toolbox for Ready-made Economic Experiments
\end{description}

% ============================================================================
% REFERENCES
% ============================================================================
\section{References}

See Master Bibliography (bcm2\_master\_references.bib) for complete citations.

\nocite{bcm_master}

\textbf{Primary Fischbacher Sources} (19 papers integrated):

\begin{itemize}
    \item Fischbacher, U., Gächter, S., \& Fehr, E. (2001). Are People Conditionally Cooperative? \textit{Econ Letters}.
    \item Fehr, E., \& Fischbacher, U. (2002). Why Social Preferences Matter. \textit{EJ}.
    \item Fehr, E., \& Fischbacher, U. (2003). The Nature of Human Altruism. \textit{Nature}.
    \item Fehr, E., \& Fischbacher, U. (2004). Third-Party Punishment. \textit{E\&HB}.
    \item Fischbacher, U. (2007). z-Tree. \textit{Experimental Economics}.
    \item Fischbacher, U., \& Gächter, S. (2010). Social Preferences, Beliefs, and Dynamics. \textit{AER}.
    \item Fischbacher, U., \& Föllmi-Heusi, F. (2013). Lies in Disguise. \textit{JEEA}.
\end{itemize}

\end{document}
